\section{Galois Theory}

\exercise{53.1}{The embedding index of a multiquadratic field}
The characteristic is zero, so the extension is separable. Thus its index, the number of embeddings fixing $\Q$, equals its given degree: $\{K:\Q\}=8$.

\exercise{53.2}{The group order of a multiquadratic field}
All independent sign changes of $\sqrt2,\sqrt3,\sqrt5$ are automorphisms. They give $2^3=8$ maps, exhausting the degree bound. Thus $|G(K/\Q)|=8$.

\exercise{53.3}{The subgroup fixing the bottom field}
By definition $\lambda(\Q)=G(K/\Q)$, so its order is 8. The whole rational Galois group fixes the bottom field.

\exercise{53.4}{Fixing two specified radicals}
The field $\Q(\sqrt2,\sqrt3)$ has degree 4 over $\Q$. The subgroup fixing it has order $[K:\Q(\sqrt2,\sqrt3)]=8/4=2$, consisting of the identity and the sign change of $\sqrt5$.

\exercise{53.5}{Fixing a product radical}
The field $\Q(\sqrt6)$ has degree 2 over $\Q$, so its fixing subgroup has order $8/2=4$. Equivalently the signs of $\sqrt2$ and $\sqrt3$ must agree, while the sign of $\sqrt5$ is arbitrary.

\exercise{53.6}{Fixing a triple-product radical}
The quadratic field $\Q(\sqrt{30})$ has degree 2 over $\Q$, so its fixing subgroup has order 4. Its sign changes are exactly the choices with a positive product of the three radical signs.

\exercise{53.7}{Fixing a sum of two radicals}
Put $\alpha=\sqrt2+\sqrt6$. Its square is $8+4\sqrt3$, so $\sqrt3=(\alpha^2-8)/4\in\Q(\alpha)$. Also $\sqrt2=\alpha/(1+\sqrt3)$ lies in this field. Thus $\Q(\alpha)=\Q(\sqrt2,\sqrt3)$ has degree 4; its fixing subgroup has order $8/4=2$.

\exercise{53.8}{Fixing the entire top field}
An automorphism fixing every element of $K$ is the identity. Therefore $\lambda(K)$ is the trivial subgroup, of order 1.

\exercise{53.9}{The group of fourth roots of unity}
The splitting field of $x^4-1$ is $\Q(i)$. The two automorphisms send $i$ to $i$ and $-i$. Their composition table is that of the cyclic group $\Z_2$; the nonidentity map is complex conjugation.

\exercise{53.10}{A relative finite-field Frobenius}
Since $729=9^3$, the degree of $\operatorname{GF}(729)$ over $\operatorname{GF}(9)$ is 3. Its Galois group is cyclic of order 3, generated by $a\mapsto a^9$. The third power is $a\mapsto a^{729}=a$; a smaller positive power cannot fix all 729 elements by the polynomial root bound.

\exercise{53.11}{The cubic splitting-field correspondence}
Let $r=\sqrt[3]2$, $\omega=(-1+i\sqrt3)/2$, $K=\Q(r,\omega)$, $\rho(r)=r\omega$, $\rho(\omega)=\omega$, and let $\kappa$ be complex conjugation.
\begin{enumerate}[label=(\alph*)]
\item The six maps and their images are
\[
\begin{array}{c|cc}
\sigma&\sigma(r)&\sigma(i\sqrt3)\\\hline
1&r&i\sqrt3\\
\rho&r\omega&i\sqrt3\\
\rho^2&r\omega^2&i\sqrt3\\
\kappa&r&-i\sqrt3\\
\rho\kappa&r\omega&-i\sqrt3\\
\rho^2\kappa&r\omega^2&-i\sqrt3
\end{array}
\]
These account for all six rational embeddings of the degree-six splitting field.
\item Their faithful action on the three cubic roots identifies the group with $S_3$. In particular $\rho^3=\kappa^2=1$ and $\kappa\rho\kappa=\rho^{-1}$.
\item The correspondence is
\[
\begin{array}{c|c}
\text{subgroup}&\text{fixed field}\\\hline
S_3&\Q\\
\langle\rho\rangle&L=\Q(\omega)\\
\langle\kappa\rangle&E_0=\Q(r)\\
\langle\rho\kappa\rangle&E_2=\Q(r\omega^2)\\
\langle\rho^2\kappa\rangle&E_1=\Q(r\omega)\\
\{1\}&K
\end{array}
\]
For instance $\rho\kappa(r\omega^2)=r\omega\cdot\omega^{-2}=r\omega^2$. The field degrees (3 for the $E_j$, 2 for $L$) and the Galois correspondence prove these are the full fixed fields. The diagrams below have only the displayed covering relations; invert one to obtain the other.
\begin{center}\setlength{\unitlength}{1pt}
\begin{picture}(400,130)
\put(200,120){\makebox(0,0){$K$}}
\put(45,65){\makebox(0,0){$E_0$}}\put(145,65){\makebox(0,0){$E_1$}}\put(255,65){\makebox(0,0){$E_2$}}\put(355,65){\makebox(0,0){$L$}}
\put(200,10){\makebox(0,0){$\Q$}}
\Line(190,113)(50,75)\Line(196,110)(148,75)\Line(204,110)(252,75)\Line(210,113)(350,75)
\Line(50,55)(190,17)\Line(148,55)(196,20)\Line(252,55)(204,20)\Line(350,55)(210,17)
\end{picture}
\begin{picture}(400,130)
\put(200,120){\makebox(0,0){$S_3$}}
\put(45,65){\makebox(0,0){$\langle\kappa\rangle$}}\put(145,65){\makebox(0,0){$\langle\rho^2\kappa\rangle$}}\put(255,65){\makebox(0,0){$\langle\rho\kappa\rangle$}}\put(355,65){\makebox(0,0){$\langle\rho\rangle$}}
\put(200,10){\makebox(0,0){$\{1\}$}}
\Line(190,113)(50,75)\Line(196,110)(148,75)\Line(204,110)(252,75)\Line(210,113)(350,75)
\Line(50,55)(190,17)\Line(148,55)(196,20)\Line(252,55)(204,20)\Line(350,55)(210,17)
\end{picture}\end{center}
\end{enumerate}

\exercise{53.12}{The group of two quadratic factors}
Since $x^4-5x^2+6=(x^2-2)(x^2-3)$, its splitting field is $\Q(\sqrt2,\sqrt3)$ of degree 4. The independent sign flips of the two radicals commute and each has order 2. Thus its group is $\Z_2\times\Z_2$, the Klein four-group.

\exercise{53.13}{The group of third roots of unity}
The splitting field is $\Q(\omega)$ with $\omega^2+\omega+1=0$. Its two maps send $\omega$ to $\omega$ and $\omega^2$, so the group is cyclic of order 2.

\exercise{53.14}{An abstract group does not determine its field}
Take $K_1=\Q(\sqrt2)$ and $K_2=\Q(\sqrt3)$ over $F=\Q$. Both have Galois group $\Z_2$. They are not isomorphic as fields: any isomorphism fixes $\Q$ and would send $\sqrt2$ to a square root of 2 in $K_2$. Writing $(a+b\sqrt3)^2=2$ forces $ab=0$, leading to $a^2=2$ or $b^2=2/3$, neither with rational solutions. Thus such an image does not exist.

\exercise{53.15}{Ten statements about the Galois correspondence}
Here ``finite normal'' has Fraleigh's definition: a \emph{separable splitting extension}, often called a finite Galois extension.
\begin{enumerate}[label=(\alph*)]
\item False. The subgroup-to-fixed-field correspondence is one-to-one.
\item False. Inclusion reverses: $E\subseteq L$ implies $\lambda(L)\subseteq\lambda(E)$.
\item True. The same splitting polynomials work over $E$, and separability descends along the tower.
\item True. Each degree equals the order of its Galois group.
\item True. A normal subgroup corresponds exactly to a normal intermediate extension of $F$.
\item False. $\Q(\zeta_5)/\Q$ is simple but has cyclic group of order 4, which has a proper subgroup of order 2.
\item False. A quadratic separable extension has simple Galois group $\Z_2$.
\item True. Finite-field Galois groups are cyclic.
\item False in characteristic 2 with this definition. $\Z_2(t^{1/2})/\Z_2(t)$ has degree 2 but is inseparable, hence is not ``finite normal'' as Fraleigh defines it (it is normal in the modern splitting-only sense).
\item True. A quadratic in characteristic different from 2 is separable and splits in its root field.
\end{enumerate}

\exercise{53.16}{Abelian Galois extensions in a tower}
The group $G(K/E)$ is a subgroup of the abelian group $G(K/F)$, hence is abelian. If $E/F$ is normal, restriction gives $G(E/F)\cong G(K/F)/G(K/E)$. A quotient of an abelian group is abelian. Thus $K/E$ and $E/F$ are both abelian extensions in the requested sense. In fact every subgroup of an abelian group is normal, so the normality of intermediate $E/F$ follows automatically here.

\exercise{53.17}{Norm and trace lie in the base field}
For any $\tau\in G(K/F)$, left multiplication $\sigma\mapsto\tau\sigma$ permutes the group. Hence
\[\tau\left(\prod_\sigma\sigma(\alpha)\right)=\prod_\sigma(\tau\sigma)(\alpha)=\prod_\sigma\sigma(\alpha),\]
and the analogous equality holds for the sum. Thus norm and trace are fixed by the entire group. Its fixed field is exactly $F$ by the Galois correspondence, so both belong to $F$.

\exercise{53.18}{Eight norm and trace calculations}
The four embeddings independently change the signs of $\sqrt2,\sqrt3$.
\begin{enumerate}[label=(\alph*)]
\item $N(\sqrt2)=(\sqrt2)^2(-\sqrt2)^2=4$.
\item $N(\sqrt2+\sqrt3)=(\sqrt2+\sqrt3)(\sqrt2-\sqrt3)(-\sqrt2+\sqrt3)(-\sqrt2-\sqrt3)=1$, since the first pair is $-1$ and the second pair is $-1$.
\item $N(\sqrt6)=(\sqrt6)^2(-\sqrt6)^2=36$.
\item Every embedding fixes 2, so $N(2)=2^4=16$.
\item The two copies each of $\sqrt2,-\sqrt2$ sum to 0.
\item The four sign choices of $\sqrt2+\sqrt3$ sum to 0.
\item The two copies each of $\sqrt6,-\sqrt6$ sum to 0.
\item The four copies of 2 sum to 8.
\end{enumerate}

\exercise{53.19}{Norm and trace from minimal-polynomial coefficients}
Because $K=F(\alpha)$ is finite normal, its $n$ automorphisms give exactly the $n$ distinct roots of the minimal polynomial of $\alpha$. Thus
\[x^n+a_{n-1}x^{n-1}+\cdots+a_0=\prod_{\sigma\in G(K/F)}(x-\sigma(\alpha)).\]
Comparing constant terms gives $a_0=(-1)^n\prod_\sigma\sigma(\alpha)$, so (a) $N_{K/F}(\alpha)=(-1)^na_0$. Comparing the coefficient of $x^{n-1}$ gives (b) $T_{K/F}(\alpha)=-a_{n-1}$.

\exercise{53.20}{The group order divides the factorial}
Let $\Omega$ be the set of distinct roots of $f$ in its splitting field $K$, and let $m=|\Omega|\le n$. Because all irreducible factors are separable, $K/F$ is finite normal in the book's sense. The faithful action of $G(K/F)$ on its generating roots embeds it into $S_m$. Lagrange's theorem gives $|G(K/F)|\mid m!$, and $m!\mid n!$. Hence the order divides $n!$. Repeated irreducible factors in $f$ cause no difficulty when one uses distinct roots.

\exercise{53.21}{The natural permutation group of a polynomial}
For $\sigma\in G(K/F)$, its values on the roots define a permutation, since applying $\sigma$ to $f(\alpha)=0$ gives $f(\sigma\alpha)=0$. Restriction respects composition. Its kernel fixes every root, hence fixes the field they generate over $F$, namely $K$, so its kernel is trivial. Thus the Galois group is naturally isomorphic to its image, a subgroup of the permutations of the distinct roots.

\exercise{53.22}{Cyclotomic extensions over a general field}
\begin{enumerate}[label=(\alph*)]
\item The powers of a primitive $n$th root $\zeta$ give all roots of $x^n-1$. Its derivative $nx^{n-1}$ has no common root with it under the characteristic hypothesis. Thus $F(\zeta)$ is a finite separable splitting extension, hence finite normal.
\item Any automorphism sends $\zeta$ to another primitive root $\zeta^a$, where $a$ is invertible modulo $n$. The map assigning $a$ to that automorphism is injective, and composition multiplies exponents modulo $n$. It embeds the group into the abelian unit group $(\Z/n\Z)^\times$. Therefore the Galois group is abelian, even when not every invertible exponent occurs.
\end{enumerate}

\exercise{53.23}{Cyclic extensions and their intermediate fields}
\begin{enumerate}[label=(\alph*)]
\item A subgroup of a cyclic group is cyclic, so $G(K/E)$ is cyclic. If $E/F$ is normal, its group is the quotient $G(K/F)/G(K/E)$, also cyclic. In fact the normality follows automatically because every subgroup of a cyclic group is normal.
\item If $|G(K/F)|=n$, a cyclic group has exactly one subgroup of each order dividing $n$. An intermediate field of degree $d$ over $F$ corresponds to a subgroup of index $d$, or order $n/d$. Hence for each divisor $d$ of $n$ there is exactly one such intermediate field.
\end{enumerate}

\exercise{53.24}{The polynomial formed from every automorphism}
\begin{enumerate}[label=(\alph*)]
\item Each $\tau\in G(K/F)$ permutes the factors of $f(x)=\prod_\sigma(x-\sigma\alpha)$, since $\sigma\mapsto\tau\sigma$ permutes the group. Thus every coefficient is fixed by the group, and belongs to $F$.
\item All distinct conjugates of $\alpha$ occur, by the extension theorem, and each occurs the same number of times: the maps taking $\alpha$ to a specified conjugate form a coset of its stabilizer $G(K/F(\alpha))$. Thus
\[f(x)=\operatorname{irr}(\alpha,F)^{[K:F(\alpha)]}.\]
Since the minimal polynomial is nonconstant, this equals the minimal polynomial exactly when the exponent is 1, namely when $K=F(\alpha)$.
\end{enumerate}

\exercise{53.25}{A compositum corresponds to an intersection}
An automorphism fixes the join $E\vee L$ if and only if it fixes both $E$ and $L$: one direction follows from containment, and the other because it then fixes every rational expression in elements of the two fields. Therefore
\[G(K/(E\vee L))=G(K/E)\cap G(K/L).\]
This is the reversal of taking the smallest field containing two given fields.

\exercise{53.26}{An intersection corresponds to a generated subgroup}
Put $H=\langle G(K/E),G(K/L)\rangle$. An element is fixed by $H$ exactly when it is fixed by both generating subgroups, so
\[K_H=K_{G(K/E)}\cap K_{G(K/L)}=E\cap L.\]
The bijective Galois correspondence now gives
\[G(K/(E\cap L))=\langle G(K/E),G(K/L)\rangle.\]
One must generate a subgroup, rather than merely take the union, which need not be a group.
